% =====================================================================
% axona-doc-preamble.tex — shared LaTeX preamble for the programmer-guide
% quartet (Quick-Start, Programmer Guide, API Reference, Services Guide).
%
% These docs are markdown-native; the PDF is built by converting the .md to
% a tufte-handout .tex with pandoc (-H this file) and compiling with tectonic.
% See render.sh. Modeled on explainer/Axona-Explainer.tex (the flat-essay
% original); the tweaks below make tufte-handout work for code/table-heavy
% reference material.
% =====================================================================

\definecolor{rust}{RGB}{185,78,54}
\definecolor{rustsoft}{RGB}{212,168,150}
\definecolor{inkmuted}{RGB}{102,102,102}
\usepackage{microtype}

% Explainer-matched page furniture: tufte-handout's letterspaced running
% heads fail on certain page breaks (the SOUL \MakeTextLowercase issue the
% explainer works around), so use the same fancyhdr style — no running
% head, a centered page number in the footer.
\usepackage{fancyhdr}
\fancypagestyle{ndhtplain}{%
  \fancyhf{}%
  \renewcommand{\headrulewidth}{0pt}%
  \fancyfoot[C]{\thepage}%
}

% Explainer-matched heading-orphan protection: reserve vertical room before
% each (sub)section heading so a heading never sits alone at a page bottom.
\usepackage{etoolbox}
\usepackage{needspace}
\pretocmd{\section}{\needspace{6\baselineskip}}{}{}
\pretocmd{\subsection}{\needspace{4\baselineskip}}{}{}

% XeTeX (tectonic) unicode monospace: renders … → and the box-drawing ASCII
% diagrams that the default CM/LM mono cannot. Menlo ships with macOS.
\setmonofont{Menlo}[Scale=0.78]

% tufte-handout ships \subsubsection as an error stub (it wants flat essays).
% A run-in fallback lets nested reference headings (§4.2, §12.3, …) compile.
\renewcommand{\subsubsection}[1]{\medskip\par\noindent{\normalsize\bfseries #1}\par\nobreak\smallskip}

% These docs use no margin notes, so reclaim tufte's wide right margin: widen
% the text block so code blocks and API tables don't overflow it.
\usepackage{geometry}
\geometry{left=1in, textwidth=6.1in, marginparsep=0.2in, marginparwidth=0.9in,
          top=1in, headsep=0.4in}

% Wrap long code lines (cheat-sheet comments run to ~130 chars) with a ↪ marker
% instead of overflowing. Applies to pandoc's fancyvrb Highlighting/Verbatim.
\usepackage{fvextra}
\fvset{breaklines=true, breakanywhere=true}

% A few technical symbols appear in prose / inline code that Latin Modern and
% Menlo don't cover; map them to math equivalents so they render (rather than
% dropping to tofu). Only maps chars actually used — leaves → alone (it renders).
\usepackage{newunicodechar}
\newunicodechar{≈}{\ensuremath{\approx}}
\newunicodechar{≠}{\ensuremath{\neq}}
\newunicodechar{≥}{\ensuremath{\geq}}
\newunicodechar{≤}{\ensuremath{\leq}}
\newunicodechar{×}{\ensuremath{\times}}
\newunicodechar{·}{\ensuremath{\cdot}}
\newunicodechar{‖}{\ensuremath{\|}}
